\section{Circuits for iterated functions}

\begin{result}
\subsection{Fixed powers of arbitrary finite-state maps}
Given a native De Morgan circuit of size $s$ computing a map $f$ on $n$-bit
states, every fixed natural exponent $T$ and integer $1\le L\le2^n$ admit a
circuit for $f^T$ of size at most
\[
40000\left(Ls+(n+1)^3\left(L+\left\lfloor\frac{2^n}{L}\right\rfloor+1\right)\right).
\]
The map need not be invertible. The bound is independent of $T$, but $T$ is fixed
when the circuit is chosen; it is not an input to this circuit. Checkpoints are
chosen nonconstructively, without a certified preprocessing runtime.
\citep{Schlesinger2026IteratedFunctions}
\decl{Complexity/Computability/Circuit/Boolean/Iteration.lean}{227}{Cslib.Circuits.Boolean.Iteration.fixed_power}
\end{result}

\begin{result}
\subsection{An input time for an arbitrary finite-state map}
With $f,s,n,L$ as in the fixed-power result, define
\[
B=L s+(n+1)^3\left(L+\left\lfloor\frac{2^n}{L}\right\rfloor+1\right).
\]
For every $m\in\mathbb N$, one circuit on $n+m$ inputs computes $f^T(x)$ for
all $n$-bit states $x$ and all $m$-bit times $0\le T<2^m$, using at most
\[
m(40000B+4n)
\]
gates. The time is supplied least-significant bit first. This theorem allows
noninvertible maps; its linear dependence on the time width differs from the
additive sparse-counter bound below. Circuit existence is nonuniform, with no
certified checkpoint-construction runtime.\citep{Schlesinger2026IteratedFunctions}
\decl{Complexity/Computability/Circuit/Boolean/Iteration/InputTime.lean}{146}{Cslib.Circuits.Boolean.Iteration.variable_time}
\end{result}

\begin{result}
\subsection{Batched iteration with a shared table}
Let a size-$s$ native De Morgan circuit compute $f$ on $n$ bits. For integers
$0\le r\le2^n$ and $1\le L\le2^n$, put
\[
B_r=rLs+(n+1)^3\left(rL+\left\lfloor\frac{2^n}{L}\right\rfloor+1\right).
\]
For every fixed $T\in\mathbb N$, one circuit applies $f^T$ to $r$ independent
states with at most $100000B_r$ gates. With a separate $m$-bit input time for
each state, a circuit of size at most
\[
m(100000B_r+4rn)
\]
computes all $r$ answers. The table term $\lfloor2^n/L\rfloor$ is charged once
in $B_r$, rather than once per state. All gates, including constants, are
charged. These are nonuniform size bounds, not certified preprocessing or
circuit-generation times.\citep{Schlesinger2026IteratedFunctions}
\decl{Complexity/Computability/Circuit/Boolean/Iteration/Batch.lean}{253}{Cslib.Circuits.Boolean.Iteration.batched_fixed_power}
\decl{Complexity/Computability/Circuit/Boolean/Iteration/Batch.lean}{395}{Cslib.Circuits.Boolean.Iteration.batched_variable_time}
\end{result}

\begin{result}
\subsection*{Sparse-counter conventions}
Write $W_n=\mathbb Z/2^n\mathbb Z$ for the $n$-bit words, represented
least-significant bit first, and let $c(x)=x+1$. For maps on $W_n$, put
\[
D(f,g)=\{x\in W_n:f(x)\ne g(x)\}.
\]
Circuit size below counts native De Morgan gates, including constants,
comparisons, selection and addition. These are nonuniform existence results:
they provide no efficient algorithm for constructing the correction table.
The port follows Schlesinger's iterated-functions development; the source credits Moreno and Rivera for the related block characterization
of almost-commuting permutations.
\citep{Schlesinger2026IteratedFunctions,MorenoRivera2016}
\end{result}

\begin{result}
\subsection{Maps commuting with an edited counter}
For any maps $f,h:W_n\to W_n$, if $h$ is injective and $h\circ f=f\circ h$, then
there is a circuit computing all $n$ output bits of $h$ with size at most
\[
100(n+1)(|D(f,c)|+1).
\]
There is no assumption that $f$ is a permutation or that a small circuit for $f$
is supplied. The assertion includes $n=0$.
\decl{Complexity/Computability/Circuit/Boolean/Iteration/SparseDisplacement.lean}{197}{Cslib.Circuits.Boolean.Iteration.sparse_counter_centralizer}
\end{result}

\begin{result}
\subsection{Fixed signed powers and coordinate transfer}
For every permutation $f$ of $W_n$ and every integer $T$, a circuit computes $f^T$
with size at most
\[
100(n+1)(|D(f,c)|+1).
\]
The size bound is independent of the magnitude and sign of $T$. The quantifiers
are ``for every $T$, there exists a circuit'': $T$ is fixed, not an additional
input to a single circuit.
\decl{Complexity/Computability/Circuit/Boolean/Iteration/SparseDisplacement.lean}{208}{Cslib.Circuits.Boolean.Iteration.sparse_counter_power}

\end{result}

\begin{result}
\textbf{Coordinate transfer.}
Let $f,g,e$ be permutations of $W_n$ such that $e(g(x))=e(x)+1$ for every $x$.
Suppose circuits $E$ and $E^{-1}$ compute $e$ and its inverse, respectively.
For every $T\in\mathbb Z$, there is a circuit for $f^T$ with size at most
\[
|E|+|E^{-1}|+100(n+1)(|D(f,g)|+1).
\]
Both coordinate conversions are charged. The theorem does not assert that an
arbitrary cycle has cheap coordinates or complete the general iterated-function
circuit-bound program.
\decl{Complexity/Computability/Circuit/Boolean/Iteration/SparseDisplacement.lean}{217}{Cslib.Circuits.Boolean.Iteration.sparse_coordinate_power}
\end{result}

\begin{result}
\subsection*{An input iteration count}
Keep the notation $W_n$, $c$ and $D$ introduced above. Here a circuit
reads an $n$-bit initial state $x$ and an $m$-bit nonnegative integer $T$, both
least-significant bit first. The permutation is fixed when choosing the circuit.
The chart and circuit are chosen nonuniformly; no efficient preprocessing or
circuit-generation algorithm is asserted.\citep{Schlesinger2026IteratedFunctions}
\end{result}

\begin{result}
\subsection{Sparse-counter input times and cycle labels}
For every permutation $f$ of $W_n$ and every $m\in\mathbb N$, there is a native
De Morgan circuit with $n+m$ inputs and $n$ outputs computing
\[
(x,T)\longmapsto f^T(x),\qquad 0\leq T<2^m,
\]
of size at most
\[
1000(n+3)(|D(f,c)|+m+1).
\]
The dependence on the edit count and time width is additive inside the
parentheses. One circuit handles all initial states and all $m$-bit times.
The bound includes constants, comparisons and selection gates, and permits
$n=0$ or $m=0$.
\decl{Complexity/Computability/Circuit/Boolean/Iteration/SparseTime.lean}{115}{Cslib.Circuits.Boolean.Iteration.sparse_counter_variable_time}

\end{result}

\begin{result}
\textbf{Canonical coordinates.}
For every permutation $f$ of $W_n$, there is a circuit with $n$ inputs and
$n+2(n+2)$ outputs computing the canonical cycle representative, cycle length
and phase of the input. Its size is at most
\[
500(n+3)(|D(f,c)|+1).
\]
The representative is the numerically least word on the cycle. The phase is
the forward distance from that representative, modulo the cycle length.
The period and phase each occupy $n+2$ bits, so even a period of $2^n$ fits.
\decl{Complexity/Computability/Circuit/Boolean/Iteration/SparseTime.lean}{126}{Cslib.Circuits.Boolean.Iteration.sparse_counter_labels}

\end{result}

\begin{result}
\textbf{Coordinate transfer.}
Let $f,g,e$ be permutations of $W_n$ satisfying $e(g(x))=e(x)+1$, with supplied
circuits $E$ and $E^{-1}$ for $e$ and its inverse. For every time width $m$, a
circuit computing $(x,T)\mapsto f^T(x)$ has size at most
\[
|E|+|E^{-1}|+1000(n+3)(|D(f,g)|+m+1).
\]
The $m$ time bits pass through unchanged to the iteration stage. Both state
coordinate conversions are charged once, independently of $m$.
\decl{Complexity/Computability/Circuit/Boolean/Iteration/SparseTime.lean}{147}{Cslib.Circuits.Boolean.Iteration.sparse_coordinate_variable_time}
\end{result}

\begin{result}
\subsection{Sparse-counter navigation, visit schedules, and hitting times}
For any self-map $f:\alpha\to\alpha$ on a finite linearly ordered state space
and states $x,y\in\alpha$, if $d$ is the first time with $f^d(x)=y$ and $p$ is
the minimal period of $y$, then $f^t(x)=y$ if and only if $t=d+kp$ for some
$k\in\mathbb N$ (with $p=0$ giving a single visit at $d$), and the number of
visits in $\{0,\dots,T\}$ equals $0$ when $T<d$, $1$ when $T\ge d$ and $p=0$,
and $\lfloor(T-d)/p\rfloor+1$ when $T\ge d$ and $p>0$.

On $W_n$, a visit-schedule record has width $1+(n+1)+(n+1)$ bits: a
reachability bit, an $(n+1)$-bit first hitting time field (which equals the
first hitting time $d<2^n$ when the reachability bit is true), and an
$(n+1)$-bit return period $p\le 2^n$. For a permutation $f$ of $W_n$, given any
circuit $L$ computing the canonical cycle label of $f$ on $W_n$, a circuit on
$2n$ inputs computes the visit-schedule record for input pairs $(x,y)$ using at
most $2|L|+200(n+3)$ native De Morgan gates. In particular, for every
permutation $f$ of $W_n$, there is a circuit on $2n$ inputs of size at most
\[
1200(n+3)(|D(f,c)|+1)
\]
computing the visit-schedule record (which characterizes all visit times
$f^t(x)=y$ and their count up to any deadline $T$), as well as a circuit with
$1+(n+1)$ outputs of the same size bound computing the reachability flag and,
when that flag is true, the first hitting time $d<2^n$. If $f,g,e$ are
permutations of $W_n$ with $e(g(x))=e(x)+1$ and $E$ computes $e$, then a
circuit of size at most
\[
2|E|+1200(n+3)(|D(f,g)|+1)
\]
computes the visit-schedule record for $f$; only the forward coordinate circuit
$E$ (evaluated on both input states) is needed, not an inverse circuit
$E^{-1}$. The bounded visit count is a mathematical function of the output
schedule bits and $T$, not an additional circuit taking $T$ as input, and chart
selection remains nonuniform.\citep{Schlesinger2026IteratedFunctions}
\declat{9233ca057c096485f3225a55bc531e7c72f7a473}{Complexity/Foundations/Combinatorics/FunctionalGraph/HittingTime.lean}{276}{Cslib.FunctionalGraph.visit_schedule_of_first}
\declat{9233ca057c096485f3225a55bc531e7c72f7a473}{Complexity/Foundations/Combinatorics/FunctionalGraph/HittingTime.lean}{305}{Cslib.FunctionalGraph.visit_schedule}
\declat{9233ca057c096485f3225a55bc531e7c72f7a473}{Complexity/Foundations/Combinatorics/FunctionalGraph/HittingTime.lean}{324}{Cslib.FunctionalGraph.arithmeticProgressionCount_spec}
\declat{9233ca057c096485f3225a55bc531e7c72f7a473}{Complexity/Foundations/Combinatorics/FunctionalGraph/HittingTime.lean}{368}{Cslib.FunctionalGraph.visit_count_of_first}
\declat{9233ca057c096485f3225a55bc531e7c72f7a473}{Complexity/Computability/Circuit/Boolean/Iteration/ResidueAdd.lean}{122}{Cslib.Circuits.Boolean.Iteration.subtractResidues}
\declat{9233ca057c096485f3225a55bc531e7c72f7a473}{Complexity/Computability/Circuit/Boolean/Iteration/ResidueAdd.lean}{126}{Cslib.Circuits.Boolean.Iteration.toNat_subtractResidues}
\declat{9233ca057c096485f3225a55bc531e7c72f7a473}{Complexity/Computability/Circuit/Boolean/Iteration/ResidueAdd.lean}{145}{Cslib.Circuits.Boolean.Iteration.synthesis_subtractResidues}
\declat{9233ca057c096485f3225a55bc531e7c72f7a473}{Complexity/Computability/Circuit/Boolean/Iteration/SparseNavigation.lean}{74}{Cslib.Circuits.Boolean.Iteration.visitScheduleContains}
\declat{9233ca057c096485f3225a55bc531e7c72f7a473}{Complexity/Computability/Circuit/Boolean/Iteration/SparseNavigation.lean}{84}{Cslib.Circuits.Boolean.Iteration.visitScheduleCount_spec}
\declat{9233ca057c096485f3225a55bc531e7c72f7a473}{Complexity/Computability/Circuit/Boolean/Iteration/SparseNavigation.lean}{175}{Cslib.Circuits.Boolean.Iteration.cycleLabelVisits_spec}
\declat{9233ca057c096485f3225a55bc531e7c72f7a473}{Complexity/Computability/Circuit/Boolean/Iteration/SparseNavigation.lean}{238}{Cslib.Circuits.Boolean.Iteration.cycle_label_visit_circuit}
\declat{9233ca057c096485f3225a55bc531e7c72f7a473}{Complexity/Computability/Circuit/Boolean/Iteration/SparseNavigation.lean}{261}{Cslib.Circuits.Boolean.Iteration.sparse_counter_visit_circuit}
\declat{9233ca057c096485f3225a55bc531e7c72f7a473}{Complexity/Computability/Circuit/Boolean/Iteration/SparseNavigation.lean}{273}{Cslib.Circuits.Boolean.Iteration.sparse_counter_visit_schedule}
\declat{9233ca057c096485f3225a55bc531e7c72f7a473}{Complexity/Computability/Circuit/Boolean/Iteration/SparseNavigation.lean}{294}{Cslib.Circuits.Boolean.Iteration.sparse_counter_hitting_time}
\declat{9233ca057c096485f3225a55bc531e7c72f7a473}{Complexity/Computability/Circuit/Boolean/Iteration/SparseNavigation.lean}{328}{Cslib.Circuits.Boolean.Iteration.sparse_coordinate_visit_schedule}
\end{result}

\begin{result}
\subsection{Autonomous controllers and affine-payload iteration}
For arbitrary types $\alpha,\beta$, a controlled map
$\mathrm{controlledMap}(g,a)(u,v)=(g(u),a(u)(v))$ has an autonomous controller
$g:\alpha\to\alpha$ and a controller-selected payload map
$a:\alpha\to\beta\to\beta$. For every $T\in\mathbb N$,
$\mathrm{controlledMap}(g,a)^T(u,v)=(g^T(u),\mathrm{controlledPayload}(g,a,T,u)(v))$,
where $\mathrm{controlledPayload}(g,a,0,u)=\mathrm{id}$ and
$\mathrm{controlledPayload}(g,a,T+1,u)=a(g^T(u))\circ\mathrm{controlledPayload}(g,a,T,u)$.
If a family $\mathcal F\subseteq\beta^\beta$ contains $\mathrm{id}$, is closed under
composition, and contains every $a(u)$, then
$\mathrm{controlledPayload}(g,a,T,u)\in\mathcal F$ for all $T$ and $u$.

On $q$ controller bits and $p$ payload bits, the block map
$\mathrm{controlledBits}(g,a):\{0,1\}^{q+p}\to\{0,1\}^{q+p}$ concatenates $g(u)$ and
$a(u)(v)$ via $\mathrm{Fin.append}$, with $T$-fold iterate
$\mathrm{controlledBits}(g,a)^T=\mathrm{controlledBits}\bigl(g^T,\mathrm{controlledPayload}(g,a,T)\bigr)$.
If every $a(u)$ has a native De Morgan circuit of size at most $K$, then
$\mathrm{controlledBits}(g,a)$ has a circuit of size at most
\[
2+2^q(K+6q+4p+1),
\]
without assuming any circuit bound for $g$. When every $a(u)$ lies in a
composition-closed family $\mathcal F$ containing $\mathrm{id}$ whose members each
have a circuit of size at most $K$, every fixed iterate
$\mathrm{controlledBits}(g,a)^T$ has a circuit of size at most
$2+2^q(K+6q+4p+1)$, and for every time width $m\in\mathbb N$ a circuit on
$(q+p)+m$ inputs computes $\mathrm{inputTimeIterate}(\mathrm{controlledBits}(g,a),m)$
with size at most
\[
m\bigl(2+2^q(K+6q+4p+1)+4(q+p)\bigr).
\]

Every affine hash $x\mapsto Ax+b$ from $n$ bits to $k$ bits over $\mathbb F_2$ has a
De Morgan circuit of size at most $k(5(n+1)+1)$, and every Boolean affine map
$a:\{0,1\}^p\to_{\mathrm{aff}}\{0,1\}^p$ agrees with the affine hash of
$\mathrm{affineMapSeed}(a)$ and has a circuit of size at most $p(5(p+1)+1)$. For
any autonomous controller $g:\{0,1\}^q\to\{0,1\}^q$ and affine payload family
$a:\{0,1\}^q\to(\{0,1\}^p\to_{\mathrm{aff}}\{0,1\}^p)$, put
\[
B_{\mathrm{aff}}(q,p)=2+2^q\bigl(p(5(p+1)+1)+6q+4p+1\bigr)\le 32\cdot 2^q(q+p^2+1).
\]
Then every fixed iterate $\mathrm{controlledBits}(g,a)^T$ has a circuit of size at
most $B_{\mathrm{aff}}(q,p)$ (and at most $32K(q+p+1)^{d+2}$ whenever
$2^q\le K(q+p+1)^d$), and an $m$-bit input time costs at most
$m\bigl(B_{\mathrm{aff}}(q,p)+4(q+p)\bigr)$ gates. The controller must be autonomous
(independent of the payload), and the circuit bounds are nonuniform, hardwiring
$g^T(u)$ and the composite payload coefficients for each $T$ without certifying a
uniform circuit-construction runtime.\citep{Schlesinger2026IteratedFunctions}
\declat{c4870e5ed91116a968cfbd2b4bf33eb8f9538cb5}{Complexity/Computability/Circuit/Boolean/Iteration/AffineController.lean}{100}{Cslib.Circuits.Boolean.Iteration.controlledMap}
\declat{c4870e5ed91116a968cfbd2b4bf33eb8f9538cb5}{Complexity/Computability/Circuit/Boolean/Iteration/AffineController.lean}{104}{Cslib.Circuits.Boolean.Iteration.controlledPayload}
\declat{c4870e5ed91116a968cfbd2b4bf33eb8f9538cb5}{Complexity/Computability/Circuit/Boolean/Iteration/AffineController.lean}{110}{Cslib.Circuits.Boolean.Iteration.controlledMap_iterate}
\declat{c4870e5ed91116a968cfbd2b4bf33eb8f9538cb5}{Complexity/Computability/Circuit/Boolean/Iteration/AffineController.lean}{119}{Cslib.Circuits.Boolean.Iteration.controlledPayload_mem}
\declat{c4870e5ed91116a968cfbd2b4bf33eb8f9538cb5}{Complexity/Computability/Circuit/Boolean/Iteration/AffineController.lean}{129}{Cslib.Circuits.Boolean.Iteration.controlledBits}
\declat{c4870e5ed91116a968cfbd2b4bf33eb8f9538cb5}{Complexity/Computability/Circuit/Boolean/Iteration/AffineController.lean}{135}{Cslib.Circuits.Boolean.Iteration.controlledBits_append}
\declat{c4870e5ed91116a968cfbd2b4bf33eb8f9538cb5}{Complexity/Computability/Circuit/Boolean/Iteration/AffineController.lean}{142}{Cslib.Circuits.Boolean.Iteration.controlledBits_iterate}
\declat{c4870e5ed91116a968cfbd2b4bf33eb8f9538cb5}{Complexity/Computability/Circuit/Boolean/Iteration/AffineController.lean}{156}{Cslib.Circuits.Boolean.Iteration.controlled_circuit}
\declat{c4870e5ed91116a968cfbd2b4bf33eb8f9538cb5}{Complexity/Computability/Circuit/Boolean/Iteration/AffineController.lean}{233}{Cslib.Circuits.Boolean.Iteration.controlled_family_power}
\declat{c4870e5ed91116a968cfbd2b4bf33eb8f9538cb5}{Complexity/Computability/Circuit/Boolean/Iteration/AffineController.lean}{247}{Cslib.Circuits.Boolean.Iteration.controlled_family_variable_time}
\declat{c4870e5ed91116a968cfbd2b4bf33eb8f9538cb5}{Complexity/Computability/Circuit/Boolean/Iteration/AffineController.lean}{263}{Cslib.Circuits.Boolean.Iteration.affineHash_circuit}
\declat{c4870e5ed91116a968cfbd2b4bf33eb8f9538cb5}{Complexity/Computability/Circuit/Boolean/Iteration/AffineController.lean}{302}{Cslib.Circuits.Boolean.Iteration.affineMapSeed}
\declat{c4870e5ed91116a968cfbd2b4bf33eb8f9538cb5}{Complexity/Computability/Circuit/Boolean/Iteration/AffineController.lean}{307}{Cslib.Circuits.Boolean.Iteration.affineHash_affineMapSeed}
\declat{c4870e5ed91116a968cfbd2b4bf33eb8f9538cb5}{Complexity/Computability/Circuit/Boolean/Iteration/AffineController.lean}{319}{Cslib.Circuits.Boolean.Iteration.affine_payload_circuit}
\declat{c4870e5ed91116a968cfbd2b4bf33eb8f9538cb5}{Complexity/Computability/Circuit/Boolean/Iteration/AffineController.lean}{326}{Cslib.Circuits.Boolean.Iteration.affineControllerBudget}
\declat{c4870e5ed91116a968cfbd2b4bf33eb8f9538cb5}{Complexity/Computability/Circuit/Boolean/Iteration/AffineController.lean}{330}{Cslib.Circuits.Boolean.Iteration.affine_controller_power}
\declat{c4870e5ed91116a968cfbd2b4bf33eb8f9538cb5}{Complexity/Computability/Circuit/Boolean/Iteration/AffineController.lean}{345}{Cslib.Circuits.Boolean.Iteration.affineControllerBudget_le}
\declat{c4870e5ed91116a968cfbd2b4bf33eb8f9538cb5}{Complexity/Computability/Circuit/Boolean/Iteration/AffineController.lean}{357}{Cslib.Circuits.Boolean.Iteration.affine_controller_polynomial}
\declat{c4870e5ed91116a968cfbd2b4bf33eb8f9538cb5}{Complexity/Computability/Circuit/Boolean/Iteration/AffineController.lean}{375}{Cslib.Circuits.Boolean.Iteration.affine_controller_variable_time}
\end{result}

\begin{result}
\subsection{Checkpoint forests, interval labels, and compressed-forest navigation}
For a finite state space $\alpha$, self-map $f:\alpha\to\alpha$, and checkpoint set $H\subseteq\alpha$,
the hull $\mathrm{checkpointHull}(f,H)=\{f^t(h):h\in H,\ t\le\mu(h)\}$ closes $H$ along transient
paths up to cycle entry, and cutting outgoing edges at cyclic states ($\mu(x)=0$) forms a directed
forest with sources $\mathrm{forestSources}(f,H)=\{x\in\mathrm{hull}:\mu(x)>0\}$. Marking all
vertices in $H$ together with roots ($\mu(x)=0$) and branching vertices ($|\mathrm{children}(x)|\ge 2$)
yields $\mathrm{forestMarks}(f,H)\subseteq\mathrm{checkpointHull}(f,H)$ of size at most $2|H|$ (since
leaves lie in $H$ and $|\mathrm{branches}|+|\mathrm{roots}|\le|\mathrm{leaves}|$), and every unmarked
hull vertex has at most one predecessor in $\mathrm{forestSources}(f,H)$. Ranking the root-to-vertex
paths $\mathrm{forestPath}(f,x)$ (the reversed forward traces
$(\mathrm{forestTrace}(f,x,\mu(x)+1)).\mathrm{reverse}$) in lexicographic order assigns intervals
$[\mathrm{low}(y),\mathrm{high}(y)]\subseteq[0,|\alpha|)$ such that
$\mathrm{low}(y)\le\mathrm{low}(x)\le\mathrm{high}(y)$ if and only if $f^t(x)=y$ for some $t\le\mu(x)$.

When every orbit either repeats within $L>0$ steps or hits $H$ in fewer than $2L$ steps, every
non-root mark $x\in\mathrm{forestMarks}(f,H)$ has a unique compressed edge $\mathrm{IsForestEdge}(f,H,x,d)$
to the next mark of length $1\le d\le 2L$, whose interior vertices are unmarked and cannot merge, so
the last-predecessor key $f^{d-1}(x)$ uniquely identifies $(x,d)$ ($\mathrm{ForestEdgeTable}$). Every hull
vertex $y$ has a unique mark anchor $a$ ($\mathrm{ForestAnchor}(f,H,L,y,a)$) reaching $y$ in at most
$2L$ steps. Given supplied prefix-mark and anchor selectors (returning a mark
$h\in\mathrm{forestMarks}(f,H)$ hit from $x$ in at most $2L$ steps whenever $H$ is hit within $2L$ steps,
and an anchor $a$ satisfying
$\mathrm{ForestAnchor}(f,H,L,y,a)$), for any transient target $y$ ($\mu(y)>0$), $y$ is reachable from
$x$ if and only if either $y$ appears in the $2L$-step prefix of $x$, or the selected prefix mark $h$
from $x$ and anchor $a$ of $y$ exist, $y$ appears in the $2L$-step verification prefix from $a$ (which
rejects off-hull targets), and $\mathrm{low}(a)\le\mathrm{low}(h)\le\mathrm{high}(a)$. These
are combinatorial decompositions; Boolean circuit synthesis of the compressed-forest queries is
separate.\citep{Schlesinger2026IteratedFunctions}
\declat{5f8f7de4f8c6f0e1188e3508b6c2a616d1d7abb0}{Complexity/Foundations/Combinatorics/FunctionalGraph/CheckpointForest.lean}{70}{Cslib.FunctionalGraph.checkpointHull}
\declat{5f8f7de4f8c6f0e1188e3508b6c2a616d1d7abb0}{Complexity/Foundations/Combinatorics/FunctionalGraph/CheckpointForest.lean}{121}{Cslib.FunctionalGraph.forestSources}
\declat{5f8f7de4f8c6f0e1188e3508b6c2a616d1d7abb0}{Complexity/Foundations/Combinatorics/FunctionalGraph/CheckpointForest.lean}{146}{Cslib.FunctionalGraph.forestMarks}
\declat{5f8f7de4f8c6f0e1188e3508b6c2a616d1d7abb0}{Complexity/Foundations/Combinatorics/FunctionalGraph/CheckpointForest.lean}{182}{Cslib.FunctionalGraph.forestLeaves_subset}
\declat{5f8f7de4f8c6f0e1188e3508b6c2a616d1d7abb0}{Complexity/Foundations/Combinatorics/FunctionalGraph/CheckpointForest.lean}{193}{Cslib.FunctionalGraph.sum_forestChildren}
\declat{5f8f7de4f8c6f0e1188e3508b6c2a616d1d7abb0}{Complexity/Foundations/Combinatorics/FunctionalGraph/CheckpointForest.lean}{203}{Cslib.FunctionalGraph.forestBranches_roots_le_leaves}
\declat{5f8f7de4f8c6f0e1188e3508b6c2a616d1d7abb0}{Complexity/Foundations/Combinatorics/FunctionalGraph/CheckpointForest.lean}{237}{Cslib.FunctionalGraph.forestMarks_card}
\declat{5f8f7de4f8c6f0e1188e3508b6c2a616d1d7abb0}{Complexity/Foundations/Combinatorics/FunctionalGraph/CheckpointForest.lean}{252}{Cslib.FunctionalGraph.forestPredecessor_unique}
\declat{5f8f7de4f8c6f0e1188e3508b6c2a616d1d7abb0}{Complexity/Foundations/Combinatorics/FunctionalGraph/ForestIntervals.lean}{81}{Cslib.FunctionalGraph.forestPath}
\declat{5f8f7de4f8c6f0e1188e3508b6c2a616d1d7abb0}{Complexity/Foundations/Combinatorics/FunctionalGraph/ForestIntervals.lean}{125}{Cslib.FunctionalGraph.forestPath_prefix}
\declat{5f8f7de4f8c6f0e1188e3508b6c2a616d1d7abb0}{Complexity/Foundations/Combinatorics/FunctionalGraph/ForestIntervals.lean}{138}{Cslib.FunctionalGraph.list_prefix_interval}
\declat{5f8f7de4f8c6f0e1188e3508b6c2a616d1d7abb0}{Complexity/Foundations/Combinatorics/FunctionalGraph/ForestIntervals.lean}{178}{Cslib.FunctionalGraph.forestRank}
\declat{5f8f7de4f8c6f0e1188e3508b6c2a616d1d7abb0}{Complexity/Foundations/Combinatorics/FunctionalGraph/ForestIntervals.lean}{204}{Cslib.FunctionalGraph.forestUpper}
\declat{5f8f7de4f8c6f0e1188e3508b6c2a616d1d7abb0}{Complexity/Foundations/Combinatorics/FunctionalGraph/ForestIntervals.lean}{208}{Cslib.FunctionalGraph.forest_interval_spec}
\declat{5f8f7de4f8c6f0e1188e3508b6c2a616d1d7abb0}{Complexity/Foundations/Combinatorics/FunctionalGraph/ForestIntervals.lean}{228}{Cslib.FunctionalGraph.exists_forest_intervals}
\declat{5f8f7de4f8c6f0e1188e3508b6c2a616d1d7abb0}{Complexity/Foundations/Combinatorics/FunctionalGraph/ForestSegments.lean}{78}{Cslib.FunctionalGraph.IsForestEdge}
\declat{5f8f7de4f8c6f0e1188e3508b6c2a616d1d7abb0}{Complexity/Foundations/Combinatorics/FunctionalGraph/ForestSegments.lean}{104}{Cslib.FunctionalGraph.exists_short_forest_edge}
\declat{5f8f7de4f8c6f0e1188e3508b6c2a616d1d7abb0}{Complexity/Foundations/Combinatorics/FunctionalGraph/ForestSegments.lean}{131}{Cslib.FunctionalGraph.forest_edges_merge}
\declat{5f8f7de4f8c6f0e1188e3508b6c2a616d1d7abb0}{Complexity/Foundations/Combinatorics/FunctionalGraph/ForestSegments.lean}{158}{Cslib.FunctionalGraph.forest_edge_last_injective}
\declat{5f8f7de4f8c6f0e1188e3508b6c2a616d1d7abb0}{Complexity/Foundations/Combinatorics/FunctionalGraph/ForestSegments.lean}{165}{Cslib.FunctionalGraph.forest_hull_partition}
\declat{5f8f7de4f8c6f0e1188e3508b6c2a616d1d7abb0}{Complexity/Foundations/Combinatorics/FunctionalGraph/ForestSegments.lean}{212}{Cslib.FunctionalGraph.forest_segment_contains_of_depth}
\declat{5f8f7de4f8c6f0e1188e3508b6c2a616d1d7abb0}{Complexity/Foundations/Combinatorics/FunctionalGraph/ForestSegments.lean}{259}{Cslib.FunctionalGraph.checkpoint_before_distant_target}
\declat{5f8f7de4f8c6f0e1188e3508b6c2a616d1d7abb0}{Complexity/Foundations/Combinatorics/FunctionalGraph/ForestNavigation.lean}{73}{Cslib.FunctionalGraph.ForestAnchor}
\declat{5f8f7de4f8c6f0e1188e3508b6c2a616d1d7abb0}{Complexity/Foundations/Combinatorics/FunctionalGraph/ForestNavigation.lean}{78}{Cslib.FunctionalGraph.exists_forest_anchor}
\declat{5f8f7de4f8c6f0e1188e3508b6c2a616d1d7abb0}{Complexity/Foundations/Combinatorics/FunctionalGraph/ForestNavigation.lean}{91}{Cslib.FunctionalGraph.forestAnchor_unique}
\declat{5f8f7de4f8c6f0e1188e3508b6c2a616d1d7abb0}{Complexity/Foundations/Combinatorics/FunctionalGraph/ForestNavigation.lean}{109}{Cslib.FunctionalGraph.forest_anchor_contains}
\declat{5f8f7de4f8c6f0e1188e3508b6c2a616d1d7abb0}{Complexity/Foundations/Combinatorics/FunctionalGraph/ForestNavigation.lean}{134}{Cslib.FunctionalGraph.forest_navigation_spec}
\declat{5f8f7de4f8c6f0e1188e3508b6c2a616d1d7abb0}{Complexity/Foundations/Combinatorics/FunctionalGraph/ForestNavigation.lean}{180}{Cslib.FunctionalGraph.ForestEdgeTable}
\declat{5f8f7de4f8c6f0e1188e3508b6c2a616d1d7abb0}{Complexity/Foundations/Combinatorics/FunctionalGraph/ForestNavigation.lean}{198}{Cslib.FunctionalGraph.exists_forest_edge_table}
\declat{5f8f7de4f8c6f0e1188e3508b6c2a616d1d7abb0}{Complexity/Foundations/Combinatorics/FunctionalGraph/ForestNavigation.lean}{222}{Cslib.FunctionalGraph.ForestEdgeTable.key_injective}
\declat{5f8f7de4f8c6f0e1188e3508b6c2a616d1d7abb0}{Complexity/Foundations/Combinatorics/FunctionalGraph/ForestNavigation.lean}{227}{Cslib.FunctionalGraph.ForestEdgeTable.anchor_iff}
\end{result}
